\chapter{Supplementary Results and Sensitivity Analyses}
\label{app:supplementary-results}

\section{Matched Two-Epoch Representation Diagnosis}
The frozen32 versus fine-tuned32 diagnosis compares equal epoch and data budgets on the complete 4,530-action REASSEMBLE cohort. It contains five matched outer folds for seed~42 and five for seed~20261023, with 75 new and 25 explicitly reused visual fits. Equal budgets do not imply equal parameter counts, computation or training time. The fixed numerical branch is shared between arms, so a visual or fusion difference does not arise from refitting a different sensor model.

Table~\ref{tab:twoepoch-supplement} gives pooled failure-positive AP for each method at both seeds. The intervals resample paired recording clusters with the models held fixed. Mean-fold AP is a separate summary that weights folds differently.

\begin{table}[htbp]
\caption{Fine-Tuned Minus Frozen Average Precision After Two Epochs}\label{tab:twoepoch-supplement}
\centering\small

\begin{tabular}{llrrl}\hline
Seed & Method & Frozen & Fine-tuned & Difference [95\% interval]\\\hline
42 & Visual &0.222708&0.230609&+0.007902 [$-0.017312$, $+0.033219$]\\
42 & Averaging &0.444720&0.440946&$-0.003774$ [$-0.014238$, $+0.007851$]\\
42 & F2 &0.439363&0.438566&$-0.000796$ [$-0.006424$, $+0.006072$]\\
20261023 & Visual &0.216799&0.200517&$-0.016282$ [$-0.039636$, $+0.009790$]\\
20261023 & Averaging &0.436495&0.428489&$-0.008006$ [$-0.018180$, $+0.003279$]\\
20261023 & F2 &0.435202&0.432380&$-0.002822$ [$-0.008285$, $+0.003602$]\\\hline
\end{tabular}
\displaynote{The conditional 95\% intervals quantify contrasts between the fitted predictors. Variation across training seeds is a separate source of uncertainty.}
\end{table}
\ev{Matched two-epoch results}{\EvidenceRef{E062}{}}

At seed~42, the small gain in visual AP coincided with lower AP for both fusion methods. At the second seed, fine-tuning reduced the point estimate for all three methods. Every interval includes zero, so the direction of an adaptation effect remains uncertain.

The diagnostic compared training duration, complementary errors, branch calibration, F2 fitting and the difference between out-of-fold and assessment scores. The available comparisons did not identify a cause for the variation across seeds. Class-conditional Brier score and log loss characterize prediction errors on actions the sensor misclassified, but cannot determine whether those images contained sufficient visual evidence.\ev{Unresolved causal diagnosis}{\EvidenceRef{E062}{}}

\section{The Same-Test 753-Action Follow-up}
The 753-action cohort is exactly original outer fold~0, not an independent external dataset. It contains 77 failures from 26 recordings within the full 4,530-action, 148-recording cohort. The correspondence is established through native recording/segment identities, timestamps and labels, with older sample hashes retained as aliases. Failure prevalence is therefore different from that of the pooled five-fold assessment.

\begin{table}[htbp]
\caption{Prediction Performance on the Same 753 Held-Out Actions}\label{tab:753-supplement}
\centering\small

\begin{tabular}{lr}\hline Model/recipe&AP\\\hline
Original frozen visual&0.232805\\
C20 full-CNN outer refit, seed42&0.215321\\
Frozen CNN100&0.193776\\
Full CNN100&0.212585\\
Static sensor, original score vector&0.327705\\
Original F2 restricted to fold0&0.367155\\
Frozen CNN100 plus F2&0.334977\\
Full CNN100 plus F2&0.349094\\
Frozen CNN100 plus averaging&0.348007\\
Full CNN100 plus averaging&0.358533\\\hline
\end{tabular}
\displaynote{Models were evaluated on identical actions with intact inputs. The training, calibration, and stacking procedures differed between generations.}
\end{table}
\ev{Matched cohort reconciliation}{\EvidenceRef{E052}{}}

Fine-tuning improves the hundred-epoch visual recipe relative to its frozen hundred-epoch control, but the original frozen recipe has a higher point AP on the same actions. Similarly, full100 F2 exceeds frozen100 F2, while the existing paired interval crosses zero. The original F2 row uses a different nested stacking protocol. Neither comparison establishes that the evolving later recipe is the best general architecture.

The selected inner job originally fitted 1,752 training actions, selected on 419 validation actions and calibrated on 686. Its original assessment contained 920 inner actions. The hundred-epoch follow-up retains those fitting roles but assesses the 753 original outer actions; the 920-action bank becomes the held-out fusion-fitting bank. In contrast, the C20 outer refit uses 2,385/581/811 training/validation/calibration actions. These role differences remain even after the assessment identities match.

\section{Aggregation, Aliases and Sensitivity Boundaries}
Across all five original outer folds, pooled AP was 0.230658 for vision, 0.429392 for the static sensor and 0.447299 for F2. The C20 full-CNN model achieved pooled AP of 0.312127. The hundred-epoch follow-up assessed only 753 of those actions and produced no predictions for the other 3,777.

Several apparent comparisons reuse identical predictions. The M3 sensor row contains the original sensor vector, and the hundred-epoch fusion study uses the standalone visual outputs. Native recording and segment keys assign the 753 outer assessment actions to 26 recordings. The 28-recording list in the selected-fit metadata refers to the inner prediction bank.

MSM varied the target definition and grouping policy and compared exposure to synthetic corruption. HANDOVER evaluated its baselines without fitting learned fusion because participant-disjoint out-of-fold predictions were unavailable. ViFailback compared the original frozen representation and later matched donor conditions. Differences in targets, prevalence and grouping limit any comparison of average precision across these datasets.\ev{Sensitivity and status inventory}{\EvidenceRef{E108}{}}
